\section{Agents: A Contract Before Code}
\label{sec:agents}

\subsection{The manifest}

An agent is a folder in a git repository: a \code{manifest.yaml} and
Python code. The manifest declares who the agent is; its tools and
functions; for each function its input and output schemas and a
permission level; what it keeps (typed records and files); what
credentials it needs, by shape, never by value; which hosts it connects
to; and its dependencies. \Cref{lst:manifest} shows the Gmail agent's
manifest, trimmed.

\begin{lstlisting}[caption={Excerpt of the Gmail agent's manifest.},label=lst:manifest,captionpos=b,float=t]
agent: {id: gmail, version: "0.3.2"}
network:
  hosts: [gmail.googleapis.com]
resources:
  secrets:
    - id: google            # connected by signing in
      oauth: {provider: google, ...}
  data:
    - id: draft
      fields:
        - {name: to, type: string, storage: keys}
        - {name: status, type: select, storage: keys,
           options: [prepared, sent, unknown, discarded]}
tools:
  - id: draft
    resources: {secrets: [google], data: [draft]}
    functions:
      - id: send
        permission_level: 3
        resources: {secrets: {google: use},
                    data: {draft: [read, update]}}
        inputs:
          type: object
          additionalProperties: false
          required: [draft_ref]
          properties:
            draft_ref: {type: string,
              x-resource: {type: data, id: draft}}
\end{lstlisting}

One validator (\file{contracts/agent\_manifest.py}) serves both the
backend and the runtime, so they cannot disagree. Among its rules:
\code{network} is required---silence is not an answer, and an agent
that connects nowhere says \code{hosts: []}; hosts are DNS names,
optionally wildcarded, a field of a declared credential
(\code{from\_secret}), or \code{any}, and literal addresses are refused;
schemas are a JSON Schema subset with a closed keyword list, no
\code{\$ref} or combinators, and bounded depth; \code{permission\_level}
is required. Flags an approver sees---\code{llm}, \code{credentials},
\code{code}, \code{watch}, \code{schedulable}---are constrained:
\code{schedulable} requires level 0 or 1 and no model use, so nothing
that runs unattended on a clock can leave the platform.

\subsection{From a repository to an approval}

\begin{enumerate}[leftmargin=1.4em]
\item An organization names a source repository. The backend fetches
  it shallowly and \emph{reads it without running anything}, showing an
  administrator each manifest, its README, its file names and sizes, and
  any validation errors.
\item An administrator approves one agent at one exact commit: the
  commit whose catalog they were shown. If the source was re-read in
  between, the approval is refused (\code{source\_moved}).
\item The backend packs that folder into a canonical archive (files
  sorted, timestamps and owners zeroed, fixed compression) named by the
  SHA-256 of its bytes, and stores it. A repository later deleted or
  rewritten changes nothing the organization runs. A version once
  approved can never be approved again with a different manifest or
  different bytes.
\item The runtime fetches the archive by digest, rejects bytes whose
  digest differs before writing anything, extracts it refusing links
  and escaping paths, re-validates the manifest against the approved
  hash, builds the environment (\cref{sec:sandbox}), and starts the
  agent once to check that it implements every declared function. Only
  then does it serve the agent.
\end{enumerate}

\noindent Each approval is named by a reference the platform mints,
\code{agt\_} plus 20 hex characters (74 random bits), per
organization. Grants,
records, and the audit trail speak only this reference; the package's
own id is translated at the manifest lookup and the worker call and
nowhere else. This is what stops one organization from naming, and so
reaching, another's agent; two organizations approving the same bytes
get two separate worker processes.

\subsection{Held to the manifest on every call}

The agent's code holds no authority of its own. When it needs
something---a credential, a record, a question to the person---it
\emph{asks} on its protocol line, and each request is answered from
the manifest entry of the function now running
(\code{ResourceAccess}, \file{ai\_runtime/execution/resources.py}).
Gmail's \code{send} declared \code{draft: [read, update]}; asked to
delete a draft or read another resource, it is refused however its code
is written. Every resource id becomes \code{<approval>\_\_<id>}, written
by the runtime from the agent actually running, so a record in another
agent's slot answers as if it did not exist. A request is answered only
if it names a call that was sent to the asking worker.

The worker protocol is newline-delimited JSON with a 2\,MiB line limit;
a malformed line ends the worker. The runtime treats the process as
hostile: lying, sending garbage or too much, or acting for a call not
its own gains nothing.

\subsection{Updates}

A changed manifest at the source is offered as an update, which is a
new approval checked like the first. When the approval names different
bytes, the next call ends the old worker and starts one from the new
code. A call still running in the old worker ends with an error and is
not retried, since nobody can know what a half-finished call changed.
